\paragraph{What the Model Reproduces.} With the full city, the configuration selected for the held-out test (z02) reproduces the two corridors delimiting the official affected area---Camino El Alba and General Blanche to the north, Quebrada Honda to the south---the reach to Av.\ Padre Hurtado, and the reported impacts near Duoc UC and Clínica Universidad de los Andes. It calculates about one meter of mud next to Clínica UC Christus, comparable in magnitude to the depth reported in its emergency department \citep{cadena3_2026}. The input volume that reproduces that reach in the model, about \SI{300000}{m^3}, agrees with the independently obtained estimate from reported rainfall, assuming saturated soil and rain falling across the entire catchment surface. Agreement between estimates obtained through different procedures supports the consistency of the order of magnitude, under their respective assumptions; it does not constitute independent validation.

\paragraph{Antecedent Rainfall, Isotherm and Volume.} According to T13, \SI{75}{mm} fell in \SI{80}{min} on 31 January, and its impact assessment was smaller than those reported for October \citep{t13_2026mallas}. The intensities reported for October are lower; the mayor reported about \SI{144}{mm} of accumulated rainfall \citep{interferencia2026} and CR2 placed the \SI{0}{\celsius} isotherm above the catchment summit \citep{cr2_2026}. This information motivates scenarios with different assumptions about saturation and the area receiving rainfall; the comparison does not establish the cause of the event or the relative contribution of those factors. Under the calculation assumptions, the upper bound on the snowmelt contribution is below \SI{5}{\percent}: no snow was detected in the cloud-free part of the catchment in the Sentinel-2 image of 21 September.

\paragraph{The Role of Drainage (Hypothetical Scenarios).} In the scenarios tested, the impact on Quebrada Honda is reproduced when part of the calculated flow is not removed by the model sink representing the Francisco Bulnes Correa culvert. With the full city, the calculated discharge at the culvert inlet reaches about \SI{65}{m^3/s}; under the assumptions tested, an assumed capacity of \SI{30}{m^3/s} leaves surface flow, and the analysis places the calculated threshold of assumed capacity for removing it entirely near \SI{65}{m^3/s}, with capture of the entire cross section in the model. These values do not describe the real structure. Official hydrographic mapping represents channel continuity along Quebrada Honda downstream of the culvert inlet; in the simulations, flow not removed by that sink continues along that alignment. The analysis constrains assumed capacities under the model assumptions; it does not determine the capacity, condition or operation of the real structure during the event and does not allow responsibility to be attributed.

\paragraph{Possible Explanations for the Simulated Extent outside the Official Affected Area.} In z02, only one third of the simulated urban mud falls within the official affected area; with porous houses (z26), the fraction rises to 0.54, and with the corrected integrated hypothesis and measured rainfall (p34b), to 0.48. The reference is the official affected area, which is a management perimeter rather than a verified mud footprint. The following representations were tested; their results describe model sensitivity and their correspondence with the event remains unproven:
\begin{enumerate}
\item \textbf{Impermeable buildings.} In the initial configuration each building is a solid block; the press sources collected in the account report mud entering homes. Entry and retention within lots through porous buildings, a representation commonly used in urban flooding \citep{schubert2012}, were tested as an alternative assumption (v26, verified with a fine time step in z26). The fraction within the official affected area rises from 0.34 to 0.54, but the depth next to Clínica UC Christus falls to \SI{0.30}{m}.
\item \textbf{Coarse load.} The images consulted, described in the theoretical framework, show stones, logs and entrained vehicles. Fine-matrix rheology does not include that additional resistance; a slightly more concentrated mud mixture ($C_v=0.43$, between the one that extends too far and the one that falls short) represented it in a simple way (v25), without reducing the extent outside the official affected area (fraction inside: 0.346 versus 0.344 in z02).
\item \textbf{Vegetation and gardens.} FABDEM removes trees and roughness outside streets is moderate; increasing it (v27) raises the fraction within the official affected area to 0.388.
\end{enumerate}
There is also a limitation of the metric: the model does not distinguish rainwater from debris-flow mud, so rainfall runoff in other neighborhoods counts as mud outside the official affected area and penalizes hypotheses with hillside rainfall.

\paragraph{Sensitivity to Terrain.} In the tests, explicit representation of smaller channels is the factor that most changes the calculated flow path: carving the linear-park channel modifies the presence of simulated flow in Quebrada Honda, the assumed southern-branch alignment modifies its arrival at upper Álvaro del Portillo, and the representation of the culvert inlet changed Quebrada Honda coverage by \SI{44}{\percent}. With \SI{30}{m} terrain, these corridors must be carved from mapping; the official hydrographic network at 1:25,000 scale proved more reliable than OpenStreetMap for the southern branch.

\paragraph{Campus and Sunken Courtyard.} The simulated overflow at Álvaro del Portillo descends toward General Blanche and does not reach the courtyard with any of the fluids tested. The variants with local runoff, a more fluid mixture and representation of the clearing and lower roadway of La Plaza Sur were completed: the micro-simulations and m36 reproduce entry along La Plaza Sur $\rightarrow$ clearing $\rightarrow$ courtyard. With a fine mesh, m36 calculates \SI{0.17}{m} in the courtyard and wets \SI{58}{\percent} of La Plaza Sur (Section~\ref{sec:res-posterior}). This result is compatible with the images consulted, but does not identify the process that occurred during the event or constitute independent validation.

\paragraph{Test Using Held-Out Observations.} The only test using held-out observations, defined before comparing them with the model, had an unfavorable outcome: the selected hypothesis reaches 19 of 49 held-out affected facilities (\SI{39}{\percent}), compared with the specified \SI{70}{\percent}, and does not outperform the main-channel-only hypothesis (Section~\ref{sec:res-ronda}). The ECI of 72/100 measures compatibility with observations used to formulate the hypotheses; this test shows that such compatibility does not transfer to the fine spatial distribution of official impacts. The model also leaves substantial parts of the official affected area without mud, as well as exceeding its extent. This may be due to terrain resolution, building representation, volume or its temporal distribution, unmodeled local inflows, or differences between the definition of impacts in the official layer and the modeled depth threshold. Building representation, local resolution, volume and the temporal distribution of rainfall were tested subsequently (Section~\ref{sec:res-posterior}); those fits do not change the outcome of the held-out test and their independent validation remains untested.

\paragraph{What Improved after the Test.} Two changes in representation explain most of the improvement following the held-out test. First, allowing mud to enter lots redistributes the mixture within the official affected area. Second, the model assumes turbid water on campus and on La Plaza Sur and thick mud in the main branch, and represents the roads and confined roadways that convey it. Both have observational support (videos, Street View, Google Earth). Even so, they were chosen with knowledge of the results, so the increase in ECI (from 72 to 76--82 in fits verified with a fine time step) measures an improved fit, not validation. At closure, the integrated hypothesis obtains 79--82: 82 with an assumed \SI{35}{mm/h} and a fine time step (z32), 81 with the pulse measured at Río Molina (p34b, corrected) and 79 with that rainfall distributed over \SI{4}{h} (v33). With this representation of retention within lots, the calculated street depth next to Clínica UC Christus falls below the depth reported inside the clinic; this comparison does not establish actual storage within lots.

\paragraph{Fire and Metrics.} Severity estimated with dNBR (23 December 2025 and 6 February 2026) yields \SI{82}{\percent} with a burn signal and \SI{10}{\percent} with moderate-high or high severity; with the latter fraction, the Gartner median is \SI{148000}{m^3}. With the assumed coefficients, the balance without fire calculates about \SI{62}{\percent} of the volume and \SI{43}{\percent} of the solids; Staley M1 yields 0.28 with fire versus 0.07 without it, as an illustrative application. The corrected runs f35b and f38b wet \SIrange{179}{193}{ha}, compared with \SI{222}{ha} in p34b (\SIrange{13}{19}{\percent} less), and calculate \SIrange{0.19}{0.23}{m} next to Clínica UC Christus. The ECI of 80--82, compared with 81, shows that this index does not distinguish these scenarios (Section~\ref{sec:res-incendio}).

\paragraph{Validity.} The available observations are qualitative: streets and facilities reported by the press and authorities, and an official perimeter delimiting the management area rather than the area with verified mud. No published discharge measurements or high-water marks were located in the sources consulted. The exploratory compatibility index summarizes four criteria specified before simulation and serves to compare hypotheses with one another; it is not a measure of predictive accuracy, and differences of two or three points should not be interpreted as conclusive evidence that one hypothesis is superior to another.


\paragraph{Confidence by Result.} Instead of reducing results to a single indicator, Table~\ref{tab:confianza} assigns each a qualitative confidence assessment at closure.
\begin{table}[htbp]\centering\small
\caption{Qualitative assessment of confidence in each result at closure, conditional on assumptions and available evidence.}\label{tab:confianza}
\begin{tabular}{p{0.42\textwidth}p{0.5\textwidth}}\toprule
Result & Confidence level\\\midrule
General simulated mud corridors & Relatively high: compatible with multiple scenarios and the official affected area\\
Need to represent channels and transitions to culverts & Robust finding of model sensitivity\\
Order of magnitude of estimated mixture volume & Conditional estimate, with a wide range (\SIrange{181000}{468000}{m^3})\\
Actual culvert capacity & Not determined; only capacities tested in hypothetical scenarios\\
Calculated depth in front of buildings & Exploratory, dependent on terrain and parameters\\
Velocities, pressures and arrival times & Calculated indicators, without instrumental validation\\
\bottomrule\end{tabular}
\end{table}
